% GENERATED FILE. Edit canonical lessons, then run npm run generate:books.
% canonical-source-hash: fnv1a64:3b0181b98f0b04f9
% canonical-lessons: LA-C475-messis, LA-C475-messor, LA-C475-falx, LA-C475-spica, LA-C475-triticum

\chapter{Harvest, Reaper, Sickle, Ear of grain, Wheat}
\label{ch:la-harvest-reaper-sickle-ear-of-grain-wheat}
\hlchaptermodality{475}

\begin{chapteropening}
\textbf{By the end of this chapter:} \emph{I can talk about the grain harvest: the harvest, the reapers, the sickle, the ears of grain and the wheat.}

Complete the last lesson of chapter 475: I can talk about the grain harvest: the harvest, the reapers, the sickle, the ears of grain and the wheat.
\end{chapteropening}

\section[messis, messis]{messis, messis — the harvest, the crop}

\label{lesson:LA-C475-messis}



\begin{quote}
Before the new one: say the Latin for a ploughman, then the Latin for to harrow.
\end{quote}

\subsection*{You'll want to know: messis, messis}
\textbf{messis, messis} — \textquotedblleft{}the harvest, the crop\textquotedblright{}.

\begin{quote}
Aestāte messis est, et omnēs in agrīs labōrant.
\end{quote}

\begin{quote}
\emph{In summer it is harvest time, and everyone works in the fields.}
\end{quote}

Third declension, feminine: \emph{messis, messis f.}, accusative often \emph{messem}. From \emph{metō} 'to reap'.

\begin{cousinweb}
\emph{messis} is built on \emph{metō} 'to reap, to mow' (EtymDB; Lewis \& Short), from PIE \emph{h\textsubscript{2}meh\textsubscript{1}-} 'to mow'. English \emph{mow} and \emph{meadow} are Germanic cousins.
\end{cousinweb}

\begin{grammarlens}[title={What you've built}]
One more word for telling, explaining and describing at B1.
\end{grammarlens}

\subsection*{Your turn}
\begin{itemize}
\raggedright
  \item \emph{Say it:} \emph{messis, messis}
  \item \emph{Say it:} \emph{messis, messis}, once more
  \item \emph{Say it:} a sentence of your own with \emph{messis, messis}
  \item \emph{Recall:} say the Latin for a ploughman, then the Latin for to harrow, then say \emph{messis, messis} again
\end{itemize}

\begin{tcolorbox}[breakable,colback=teal!4,colframe=teal!35!black,title={Before you move on}]
What does messis, messis mean? (\textquotedblleft{}The harvest, the crop\textquotedblright{}.) Say it once more.
\end{tcolorbox}
\section[messor, messōris]{messor, messōris — a reaper, a harvester}

\label{lesson:LA-C475-messor}



\begin{quote}
Before the new one: say the Latin for to harrow, then the Latin for the harvest.
\end{quote}

\subsection*{You'll want to know: messor, messōris}
\textbf{messor, messōris} — \textquotedblleft{}a reaper, a harvester\textquotedblright{}.

\begin{quote}
Messōrēs sub sōle ardentī frūmentum metunt.
\end{quote}

\begin{quote}
\emph{The reapers cut the grain under the burning sun.}
\end{quote}

Third declension, masculine: \emph{messor, messōris m.}, from \emph{metō} with the agent ending \emph{-or}.

\begin{cousinweb}
\emph{messor} is built on \emph{metō} 'to reap' with the agent ending \emph{-or}, on \emph{mess-}, the stem of its supine (Lewis \& Short).
\end{cousinweb}

\begin{grammarlens}[title={What you've built}]
One more word for telling, explaining and describing at B1.
\end{grammarlens}

\subsection*{Your turn}
\begin{itemize}
\raggedright
  \item \emph{Say it:} \emph{messor, messōris}
  \item \emph{Say it:} \emph{messor, messōris}, once more
  \item \emph{Say it:} a sentence of your own with \emph{messor, messōris}
  \item \emph{Recall:} say the Latin for to harrow, then the Latin for the harvest, then say \emph{messor, messōris} again
\end{itemize}

\begin{tcolorbox}[breakable,colback=teal!4,colframe=teal!35!black,title={Before you move on}]
What does messor, messōris mean? (\textquotedblleft{}A reaper, a harvester\textquotedblright{}.) Say it once more.
\end{tcolorbox}
\section[falx, falcis]{falx, falcis — a sickle, a pruning hook}

\label{lesson:LA-C475-falx}



\begin{quote}
Before the new one: say the Latin for the harvest, then the Latin for a reaper.
\end{quote}

\subsection*{You'll want to know: falx, falcis}
\textbf{falx, falcis} — \textquotedblleft{}a sickle, a pruning hook\textquotedblright{}.

\begin{quote}
Messor falce frūmentum metit.
\end{quote}

\begin{quote}
\emph{The reaper cuts the grain with a sickle.}
\end{quote}

Third declension, feminine: \emph{falx, falcis f.} A curved blade for grain or for pruning vines.

\begin{cousinweb}
\emph{falx} is inherited from PIE \emph{d\textsuperscript{h}alk-} 'sickle, cutting tool' (EtymDB); Lewis \& Short suggest a link with \emph{flectō} 'to bend', for its curve, but that is not settled.
\end{cousinweb}

\begin{grammarlens}[title={What you've built}]
One more word for telling, explaining and describing at B1.
\end{grammarlens}

\subsection*{Your turn}
\begin{itemize}
\raggedright
  \item \emph{Say it:} \emph{falx, falcis}
  \item \emph{Say it:} \emph{falx, falcis}, once more
  \item \emph{Say it:} a sentence of your own with \emph{falx, falcis}
  \item \emph{Recall:} say the Latin for the harvest, then the Latin for a reaper, then say \emph{falx, falcis} again
\end{itemize}

\begin{tcolorbox}[breakable,colback=teal!4,colframe=teal!35!black,title={Before you move on}]
What does falx, falcis mean? (\textquotedblleft{}A sickle, a pruning hook\textquotedblright{}.) Say it once more.
\end{tcolorbox}
\section[spīca, spīcae]{spīca, spīcae — an ear of grain, a spike}

\label{lesson:LA-C475-spica}



\begin{quote}
Before the new one: say the Latin for a reaper, then the Latin for a sickle.
\end{quote}

\subsection*{You'll want to know: spīca, spīcae}
\textbf{spīca, spīcae} — \textquotedblleft{}an ear of grain, a spike\textquotedblright{}.

\begin{quote}
Spīcae plēnae et flāvae sunt; messis bona erit.
\end{quote}

\begin{quote}
\emph{The ears of grain are full and golden; the harvest will be good.}
\end{quote}

First declension, feminine: \emph{spīca, spīcae f.} The head of wheat or barley.

\begin{cousinweb}
\emph{spīca} is usually traced to a root meaning 'to be sharp, pointed' (Lewis \& Short), the root of \emph{spīna} 'thorn'. EtymDB lists it without a parent.
\end{cousinweb}

\begin{grammarlens}[title={What you've built}]
One more word for telling, explaining and describing at B1.
\end{grammarlens}

\subsection*{Your turn}
\begin{itemize}
\raggedright
  \item \emph{Say it:} \emph{spīca, spīcae}
  \item \emph{Say it:} \emph{spīca, spīcae}, once more
  \item \emph{Say it:} a sentence of your own with \emph{spīca, spīcae}
  \item \emph{Recall:} say the Latin for a reaper, then the Latin for a sickle, then say \emph{spīca, spīcae} again
\end{itemize}

\begin{tcolorbox}[breakable,colback=teal!4,colframe=teal!35!black,title={Before you move on}]
What does spīca, spīcae mean? (\textquotedblleft{}An ear of grain, a spike\textquotedblright{}.) Say it once more.
\end{tcolorbox}
\section[trīticum, trīticī]{trīticum, trīticī — wheat}

\label{lesson:LA-C475-triticum}



\begin{quote}
Before the new one: say the Latin for a sickle, then the Latin for an ear of grain.
\end{quote}

\subsection*{You'll want to know: trīticum, trīticī}
\textbf{trīticum, trīticī} — \textquotedblleft{}wheat\textquotedblright{}.

\begin{quote}
Agricola trīticum in agrō serit.
\end{quote}

\begin{quote}
\emph{The farmer sows wheat in the field.}
\end{quote}

Second declension, neuter: \emph{trīticum, trīticī n.} The best grain for bread.

\begin{cousinweb}
\emph{trīticum} is built on \emph{trītus}, the participle of \emph{terō} 'to rub, to thresh' (EtymDB; Lewis \& Short, after an ancient explanation): the grain that is threshed.
\end{cousinweb}

\begin{grammarlens}[title={What you've built}]
One more word for telling, explaining and describing at B1.
\end{grammarlens}

\subsection*{Your turn}
\begin{itemize}
\raggedright
  \item \emph{Say it:} \emph{trīticum, trīticī}
  \item \emph{Say it:} \emph{trīticum, trīticī}, once more
  \item \emph{Say it:} a sentence of your own with \emph{trīticum, trīticī}
  \item \emph{Recall:} say the Latin for a sickle, then the Latin for an ear of grain, then say \emph{trīticum, trīticī} again
\end{itemize}

\begin{tcolorbox}[breakable,colback=teal!4,colframe=teal!35!black,title={Before you move on}]
What does trīticum, trīticī mean? (\textquotedblleft{}Wheat\textquotedblright{}.) Say it once more.
\end{tcolorbox}
